\IfFileExists{stacks-project.cls}{%
\documentclass{stacks-project}
}{%
\documentclass{amsart}
}

% For dealing with references we use the comment environment
\usepackage{verbatim}
\newenvironment{reference}{\comment}{\endcomment}
%\newenvironment{reference}{}{}
\newenvironment{slogan}{\comment}{\endcomment}
\newenvironment{history}{\comment}{\endcomment}

% For commutative diagrams we use Xy-pic
\usepackage[all]{xy}

% We use 2cell for 2-commutative diagrams.
\xyoption{2cell}
\UseAllTwocells

% We use multicol for the list of chapters between chapters
\usepackage{multicol}

% This is generally recommended for better output
\usepackage{lmodern}
\usepackage[T1]{fontenc}

% For cross-file-references
\usepackage{xr-hyper}

% Package for hypertext links:
\usepackage{hyperref}

% For any local file, say "hello.tex" you want to link to please
% use \externaldocument[hello-]{hello}
\externaldocument[introduction-]{introduction}
\externaldocument[conventions-]{conventions}
\externaldocument[sets-]{sets}
\externaldocument[categories-]{categories}
\externaldocument[topology-]{topology}
\externaldocument[sheaves-]{sheaves}
\externaldocument[sites-]{sites}
\externaldocument[stacks-]{stacks}
\externaldocument[fields-]{fields}
\externaldocument[algebra-]{algebra}
\externaldocument[brauer-]{brauer}
\externaldocument[homology-]{homology}
\externaldocument[derived-]{derived}
\externaldocument[simplicial-]{simplicial}
\externaldocument[more-algebra-]{more-algebra}
\externaldocument[smoothing-]{smoothing}
\externaldocument[modules-]{modules}
\externaldocument[sites-modules-]{sites-modules}
\externaldocument[injectives-]{injectives}
\externaldocument[cohomology-]{cohomology}
\externaldocument[sites-cohomology-]{sites-cohomology}
\externaldocument[dga-]{dga}
\externaldocument[dpa-]{dpa}
\externaldocument[sdga-]{sdga}
\externaldocument[hypercovering-]{hypercovering}
\externaldocument[schemes-]{schemes}
\externaldocument[constructions-]{constructions}
\externaldocument[properties-]{properties}
\externaldocument[morphisms-]{morphisms}
\externaldocument[coherent-]{coherent}
\externaldocument[divisors-]{divisors}
\externaldocument[limits-]{limits}
\externaldocument[varieties-]{varieties}
\externaldocument[topologies-]{topologies}
\externaldocument[descent-]{descent}
\externaldocument[perfect-]{perfect}
\externaldocument[more-morphisms-]{more-morphisms}
\externaldocument[flat-]{flat}
\externaldocument[groupoids-]{groupoids}
\externaldocument[more-groupoids-]{more-groupoids}
\externaldocument[etale-]{etale}
\externaldocument[chow-]{chow}
\externaldocument[intersection-]{intersection}
\externaldocument[pic-]{pic}
\externaldocument[weil-]{weil}
\externaldocument[adequate-]{adequate}
\externaldocument[dualizing-]{dualizing}
\externaldocument[duality-]{duality}
\externaldocument[discriminant-]{discriminant}
\externaldocument[derham-]{derham}
\externaldocument[local-cohomology-]{local-cohomology}
\externaldocument[algebraization-]{algebraization}
\externaldocument[curves-]{curves}
\externaldocument[resolve-]{resolve}
\externaldocument[models-]{models}
\externaldocument[functors-]{functors}
\externaldocument[equiv-]{equiv}
\externaldocument[pione-]{pione}
\externaldocument[etale-cohomology-]{etale-cohomology}
\externaldocument[proetale-]{proetale}
\externaldocument[relative-cycles-]{relative-cycles}
\externaldocument[more-etale-]{more-etale}
\externaldocument[trace-]{trace}
\externaldocument[crystalline-]{crystalline}
\externaldocument[spaces-]{spaces}
\externaldocument[spaces-properties-]{spaces-properties}
\externaldocument[spaces-morphisms-]{spaces-morphisms}
\externaldocument[decent-spaces-]{decent-spaces}
\externaldocument[spaces-cohomology-]{spaces-cohomology}
\externaldocument[spaces-limits-]{spaces-limits}
\externaldocument[spaces-divisors-]{spaces-divisors}
\externaldocument[spaces-over-fields-]{spaces-over-fields}
\externaldocument[spaces-topologies-]{spaces-topologies}
\externaldocument[spaces-descent-]{spaces-descent}
\externaldocument[spaces-perfect-]{spaces-perfect}
\externaldocument[spaces-more-morphisms-]{spaces-more-morphisms}
\externaldocument[spaces-flat-]{spaces-flat}
\externaldocument[spaces-groupoids-]{spaces-groupoids}
\externaldocument[spaces-more-groupoids-]{spaces-more-groupoids}
\externaldocument[bootstrap-]{bootstrap}
\externaldocument[spaces-pushouts-]{spaces-pushouts}
\externaldocument[spaces-chow-]{spaces-chow}
\externaldocument[groupoids-quotients-]{groupoids-quotients}
\externaldocument[spaces-more-cohomology-]{spaces-more-cohomology}
\externaldocument[spaces-simplicial-]{spaces-simplicial}
\externaldocument[spaces-duality-]{spaces-duality}
\externaldocument[formal-spaces-]{formal-spaces}
\externaldocument[restricted-]{restricted}
\externaldocument[spaces-resolve-]{spaces-resolve}
\externaldocument[formal-defos-]{formal-defos}
\externaldocument[defos-]{defos}
\externaldocument[cotangent-]{cotangent}
\externaldocument[examples-defos-]{examples-defos}
\externaldocument[algebraic-]{algebraic}
\externaldocument[examples-stacks-]{examples-stacks}
\externaldocument[stacks-sheaves-]{stacks-sheaves}
\externaldocument[criteria-]{criteria}
\externaldocument[artin-]{artin}
\externaldocument[quot-]{quot}
\externaldocument[stacks-properties-]{stacks-properties}
\externaldocument[stacks-morphisms-]{stacks-morphisms}
\externaldocument[stacks-limits-]{stacks-limits}
\externaldocument[stacks-cohomology-]{stacks-cohomology}
\externaldocument[stacks-perfect-]{stacks-perfect}
\externaldocument[stacks-introduction-]{stacks-introduction}
\externaldocument[stacks-more-morphisms-]{stacks-more-morphisms}
\externaldocument[stacks-geometry-]{stacks-geometry}
\externaldocument[moduli-]{moduli}
\externaldocument[moduli-curves-]{moduli-curves}
\externaldocument[examples-]{examples}
\externaldocument[exercises-]{exercises}
\externaldocument[guide-]{guide}
\externaldocument[desirables-]{desirables}
\externaldocument[coding-]{coding}
\externaldocument[obsolete-]{obsolete}
\externaldocument[fdl-]{fdl}
\externaldocument[index-]{index}

% Theorem environments.
%
\theoremstyle{plain}
\newtheorem{theorem}[subsection]{Theorem}
\newtheorem{proposition}[subsection]{Proposition}
\newtheorem{lemma}[subsection]{Lemma}

\theoremstyle{definition}
\newtheorem{definition}[subsection]{Definition}
\newtheorem{example}[subsection]{Example}
\newtheorem{exercise}[subsection]{Exercise}
\newtheorem{situation}[subsection]{Situation}

\theoremstyle{remark}
\newtheorem{remark}[subsection]{Remark}
\newtheorem{remarks}[subsection]{Remarks}

\numberwithin{equation}{subsection}

% Macros
%
\def\lim{\mathop{\mathrm{lim}}\nolimits}
\def\colim{\mathop{\mathrm{colim}}\nolimits}
\def\Spec{\mathop{\mathrm{Spec}}}
\def\Hom{\mathop{\mathrm{Hom}}\nolimits}
\def\Ext{\mathop{\mathrm{Ext}}\nolimits}
\def\SheafHom{\mathop{\mathcal{H}\!\mathit{om}}\nolimits}
\def\SheafExt{\mathop{\mathcal{E}\!\mathit{xt}}\nolimits}
\def\Sch{\mathit{Sch}}
\def\Mor{\mathop{\mathrm{Mor}}\nolimits}
\def\Ob{\mathop{\mathrm{Ob}}\nolimits}
\def\Sh{\mathop{\mathit{Sh}}\nolimits}
\def\NL{\mathop{N\!L}\nolimits}
\def\CH{\mathop{\mathrm{CH}}\nolimits}
\def\proetale{{pro\text{-}\acute{e}tale}}
\def\etale{{\acute{e}tale}}
\def\QCoh{\mathit{QCoh}}
\def\Ker{\mathop{\mathrm{Ker}}}
\def\Im{\mathop{\mathrm{Im}}}
\def\Coker{\mathop{\mathrm{Coker}}}
\def\Coim{\mathop{\mathrm{Coim}}}

% Boxtimes
%
\DeclareMathSymbol{\boxtimes}{\mathbin}{AMSa}{"02}

%
% Macros for moduli stacks/spaces
%
\def\QCohstack{\mathcal{QC}\!\mathit{oh}}
\def\Cohstack{\mathcal{C}\!\mathit{oh}}
\def\Spacesstack{\mathcal{S}\!\mathit{paces}}
\def\Quotfunctor{\mathrm{Quot}}
\def\Hilbfunctor{\mathrm{Hilb}}
\def\Curvesstack{\mathcal{C}\!\mathit{urves}}
\def\Polarizedstack{\mathcal{P}\!\mathit{olarized}}
\def\Complexesstack{\mathcal{C}\!\mathit{omplexes}}
% \Pic is the operator that assigns to X its picard group, usage \Pic(X)
% \Picardstack_{X/B} denotes the Picard stack of X over B
% \Picardfunctor_{X/B} denotes the Picard functor of X over B
\def\Pic{\mathop{\mathrm{Pic}}\nolimits}
\def\Picardstack{\mathcal{P}\!\mathit{ic}}
\def\Picardfunctor{\mathrm{Pic}}
\def\Deformationcategory{\mathcal{D}\!\mathit{ef}}
\usepackage{cleveref}
\crefname{originalsection}{original section}{original sections}
\Crefname{originalsection}{Original section}{Original sections}
\makeatletter
\newlabel{bootstrap-section-applications}{{04SJ}{source}{}{sourcepointer-1}{}}
\newlabel{bootstrap-section-applications@cref}{{[originalsection][0][]04SJ}{[1][0][]source}{}{}{}}
\makeatother
\crefname{originallemma}{original lemma}{original lemmas}
\Crefname{originallemma}{Original lemma}{Original lemmas}
\makeatletter
\newlabel{examples-stacks-lemma-stack-of-finite-type-spaces}{{04UD}{source}{}{sourcepointer-2}{}}
\newlabel{examples-stacks-lemma-stack-of-finite-type-spaces@cref}{{[originallemma][0][]04UD}{[1][0][]source}{}{}{}}
\makeatother
\crefname{originalsection}{original section}{original sections}
\Crefname{originalsection}{Original section}{Original sections}
\makeatletter
\newlabel{examples-stacks-section-stack-in-groupoids-of-finite-type-spaces}{{04UH}{source}{}{sourcepointer-3}{}}
\newlabel{examples-stacks-section-stack-in-groupoids-of-finite-type-spaces@cref}{{[originalsection][0][]04UH}{[1][0][]source}{}{}{}}
\makeatother
\crefname{originaldefinition}{original definition}{original definitions}
\Crefname{originaldefinition}{Original definition}{Original definitions}
\makeatletter
\newlabel{stacks-definition-stack-in-groupoids}{{02ZI}{source}{}{sourcepointer-4}{}}
\newlabel{stacks-definition-stack-in-groupoids@cref}{{[originaldefinition][0][]02ZI}{[1][0][]source}{}{}{}}
\makeatother
\crefname{originallemma}{original lemma}{original lemmas}
\Crefname{originallemma}{Original lemma}{Original lemmas}
\makeatletter
\newlabel{lemma-spaces-fibred-in-groupoids}{{0D1E}{source}{}{sourcepointer-5}{}}
\newlabel{lemma-spaces-fibred-in-groupoids@cref}{{[originallemma][0][]0D1E}{[1][0][]source}{}{}{}}
\makeatother
\crefname{originallemma}{original lemma}{original lemmas}
\Crefname{originallemma}{Original lemma}{Original lemmas}
\makeatletter
\newlabel{lemma-spaces-diagonal}{{0D1F}{source}{}{sourcepointer-6}{}}
\newlabel{lemma-spaces-diagonal@cref}{{[originallemma][0][]0D1F}{[1][0][]source}{}{}{}}
\makeatother
\crefname{originallemma}{original lemma}{original lemmas}
\Crefname{originallemma}{Original lemma}{Original lemmas}
\makeatletter
\newlabel{examples-stacks-lemma-pre-stack-of-spaces}{{04UA}{source}{}{sourcepointer-7}{}}
\newlabel{examples-stacks-lemma-pre-stack-of-spaces@cref}{{[originallemma][0][]04UA}{[1][0][]source}{}{}{}}
\makeatother
\crefname{originallemma}{original lemma}{original lemmas}
\Crefname{originallemma}{Original lemma}{Original lemmas}
\makeatletter
\newlabel{bootstrap-lemma-descend-algebraic-space}{{0ADV}{source}{}{sourcepointer-8}{}}
\newlabel{bootstrap-lemma-descend-algebraic-space@cref}{{[originallemma][0][]0ADV}{[1][0][]source}{}{}{}}
\makeatother
\crefname{originallemma}{original lemma}{original lemmas}
\Crefname{originallemma}{Original lemma}{Original lemmas}
\makeatletter
\newlabel{spaces-descent-lemma-descending-property-finite-type}{{041U}{source}{}{sourcepointer-9}{}}
\newlabel{spaces-descent-lemma-descending-property-finite-type@cref}{{[originallemma][0][]041U}{[1][0][]source}{}{}{}}
\makeatother
\crefname{originallemma}{original lemma}{original lemmas}
\Crefname{originallemma}{Original lemma}{Original lemmas}
\makeatletter
\newlabel{spaces-descent-lemma-descending-property-finite-presentation}{{041V}{source}{}{sourcepointer-10}{}}
\newlabel{spaces-descent-lemma-descending-property-finite-presentation@cref}{{[originallemma][0][]041V}{[1][0][]source}{}{}{}}
\makeatother
\crefname{originallemma}{original lemma}{original lemmas}
\Crefname{originallemma}{Original lemma}{Original lemmas}
\makeatletter
\newlabel{spaces-descent-lemma-descending-property-flat}{{041W}{source}{}{sourcepointer-11}{}}
\newlabel{spaces-descent-lemma-descending-property-flat@cref}{{[originallemma][0][]041W}{[1][0][]source}{}{}{}}
\makeatother
\crefname{originallemma}{original lemma}{original lemmas}
\Crefname{originallemma}{Original lemma}{Original lemmas}
\makeatletter
\newlabel{spaces-descent-lemma-descending-property-proper}{{0422}{source}{}{sourcepointer-12}{}}
\newlabel{spaces-descent-lemma-descending-property-proper@cref}{{[originallemma][0][]0422}{[1][0][]source}{}{}{}}
\makeatother
\title{The stack of finite type spaces is a stack in groupoids}
\author{The Stacks Project Authors}
\date{stacks-project:quot.tex:local-snapshot-20261012; GNU FDL 1.2 or later}
\begin{document}
\maketitle
\setcounter{subsection}{1}
\noindent\textit{Editorial note.} All delivered mathematical text is the author's own, copied from the pinned Stacks Project chapter file quot.tex; no statement or proof step was rewritten, reordered or removed. The counted span is the author environment and its own proof environment. References inside the span are handed to the assembler as navigation stubs that carry the canonical tag of the cited tag only; whether each of them is load-bearing is recorded on the stub itself.\par\medskip
\noindent\textit{Allowed standard prerequisites (editorial).}\begin{itemize}
\item Slogan bootstrap-section-applications (canonical tag 04SJ). cited at proof lines 3575-3576. Any fppf descent datum for algebraic spaces is effective.
\item Lemma examples-stacks-lemma-stack-of-finite-type-spaces (canonical tag 04UD). cited at proof line 3579. The finite type case of condition (3).
\item Section examples-stacks-section-stack-in-groupoids-of-finite-type-spaces (canonical tag 04UH). cited at proof line 3581. Where the finite type case was established.
\item Definition stacks-definition-stack-in-groupoids (canonical tag 02ZI). cited at proof line 3584. Conditions (1), (2), (3) being verified.
\item Lemma lemma-spaces-fibred-in-groupoids (canonical tag 0D1E). cited at proof line 3587. Condition (1): fibred in groupoids.
\item Lemma lemma-spaces-diagonal (canonical tag 0D1F). cited at proof line 3591. Condition (2) in the fp,flat,proper case.
\item Lemma examples-stacks-lemma-pre-stack-of-spaces (canonical tag 04UA). cited at proof lines 3594-3595. Condition (2) in the finite type case.
\item Lemma bootstrap-lemma-descend-algebraic-space (canonical tag 0ADV). cited at proof line 3612. Condition (3): the descended algebraic space exists (part (2)).
\item Lemma spaces-descent-lemma-descending-property-finite-type (canonical tag 041U). cited at proof line 3614. X -> U is of finite type.
\item Lemma spaces-descent-lemma-descending-property-finite-presentation (canonical tag 041V). cited at proof line 3619. fp,flat,proper case: finite presentation descends.
\item Lemma spaces-descent-lemma-descending-property-flat (canonical tag 041W). cited at proof line 3620. fp,flat,proper case: flatness descends.
\item Lemma spaces-descent-lemma-descending-property-proper (canonical tag 0422). cited at proof line 3621. fp,flat,proper case: properness descends.
\end{itemize}\medskip
\ifdefined\hypertarget\hypertarget{sourcepointer-1}{}\fi
\noindent\textit{Source reference.} Original section 04SJ: sources/stacks-chow/tags.txt, source lines 6222--6222. Cited at proof lines 3575-3576: the slogan underlying the whole lemma; load-bearing.\par\medskip
\ifdefined\hypertarget\hypertarget{sourcepointer-2}{}\fi
\noindent\textit{Source reference.} Original lemma 04UD: sources/stacks-chow/tags.txt, source lines 6286--6286. Cited at proof line 3579: the primary reference for condition (3) in the finite type case; load-bearing.\par\medskip
\ifdefined\hypertarget\hypertarget{sourcepointer-3}{}\fi
\noindent\textit{Source reference.} Original section 04UH: sources/stacks-chow/tags.txt, source lines 6290--6290. Cited at proof line 3581: where the lemma was previously seen; load-bearing.\par\medskip
\ifdefined\hypertarget\hypertarget{sourcepointer-4}{}\fi
\noindent\textit{Source reference.} Original definition 02ZI: sources/stacks-chow/tags.txt, source lines 3913--3913. Cited at proof line 3584: the three conditions being checked; load-bearing.\par\medskip
\ifdefined\hypertarget\hypertarget{sourcepointer-5}{}\fi
\noindent\textit{Source reference.} Original lemma 0D1E: sources/stacks-chow/tags.txt, source lines 16344--16344. Cited at proof line 3587: condition (1); load-bearing.\par\medskip
\ifdefined\hypertarget\hypertarget{sourcepointer-6}{}\fi
\noindent\textit{Source reference.} Original lemma 0D1F: sources/stacks-chow/tags.txt, source lines 16345--16345. Cited at proof line 3591: condition (2) in the fp,flat,proper case; load-bearing.\par\medskip
\ifdefined\hypertarget\hypertarget{sourcepointer-7}{}\fi
\noindent\textit{Source reference.} Original lemma 04UA: sources/stacks-chow/tags.txt, source lines 6283--6283. Cited at proof lines 3594-3595: condition (2) in the finite type case; load-bearing.\par\medskip
\ifdefined\hypertarget\hypertarget{sourcepointer-8}{}\fi
\noindent\textit{Source reference.} Original lemma 0ADV: sources/stacks-chow/tags.txt, source lines 13103--13103. Cited at proof line 3612: condition (3), the descended space X exists; load-bearing.\par\medskip
\ifdefined\hypertarget\hypertarget{sourcepointer-9}{}\fi
\noindent\textit{Source reference.} Original lemma 041U: sources/stacks-chow/tags.txt, source lines 5305--5305. Cited at proof line 3614: X -> U is of finite type; load-bearing.\par\medskip
\ifdefined\hypertarget\hypertarget{sourcepointer-10}{}\fi
\noindent\textit{Source reference.} Original lemma 041V: sources/stacks-chow/tags.txt, source lines 5306--5306. Cited at proof line 3619: fp,flat,proper case; load-bearing.\par\medskip
\ifdefined\hypertarget\hypertarget{sourcepointer-11}{}\fi
\noindent\textit{Source reference.} Original lemma 041W: sources/stacks-chow/tags.txt, source lines 5307--5307. Cited at proof line 3620: fp,flat,proper case; load-bearing.\par\medskip
\ifdefined\hypertarget\hypertarget{sourcepointer-12}{}\fi
\noindent\textit{Source reference.} Original lemma 0422: sources/stacks-chow/tags.txt, source lines 5313--5313. Cited at proof line 3621: fp,flat,proper case; load-bearing.\par\medskip
\begin{lemma}
\label{lemma-spaces-stack}
The category $\Spacesstack'_{ft}$ is a stack in groupoids
over $\Sch_{fppf}$. The same is true for
$\Spacesstack'_{fp, flat, proper}$.
\end{lemma}

\begin{proof}
The reason this lemma holds is the slogan: any fppf descent datum for algebraic
spaces is effective, see Bootstrap, Section
\ref{bootstrap-section-applications}.
More precisely, the lemma for $\Spacesstack'_{ft}$ follows from
Examples of Stacks, Lemma
\ref{examples-stacks-lemma-stack-of-finite-type-spaces}
as we saw in Examples of Stacks, Section
\ref{examples-stacks-section-stack-in-groupoids-of-finite-type-spaces}.
However, let us review the proof. We need to check conditions
(1), (2), and (3) of Stacks, Definition
\ref{stacks-definition-stack-in-groupoids}.

\medskip\noindent
Property (1) we have seen in Lemma \ref{lemma-spaces-fibred-in-groupoids}.

\medskip\noindent
Property (2) follows from
Lemma \ref{lemma-spaces-diagonal} in the case of
$\Spacesstack'_{fp, flat, proper}$.
In the case of $\Spacesstack'_{ft}$ it follows
from Examples of Stacks, Lemma
\ref{examples-stacks-lemma-pre-stack-of-spaces}
(and this is really the ``correct'' reference).

\medskip\noindent
Condition (3) for $\Spacesstack'_{ft}$ is checked as follows. Suppose given
\begin{enumerate}
\item an fppf covering $\{U_i \to U\}_{i \in I}$ in $\Sch_{fppf}$,
\item for each $i \in I$ an algebraic space $X_i$ of finite type over
$U_i$, and
\item for each $i, j \in I$ an isomorphism
$\varphi_{ij} : X_i \times_U U_j \to U_i \times_U X_j$ of algebraic spaces
over $U_i \times_U U_j$ satisfying the cocycle condition over
$U_i \times_U U_j \times_U U_k$.
\end{enumerate}
We have to show there exists an algebraic space $X$ of finite type over $U$
and isomorphisms $X_{U_i} \cong X_i$ over $U_i$ recovering the
isomorphisms $\varphi_{ij}$. This follows from
Bootstrap, Lemma \ref{bootstrap-lemma-descend-algebraic-space} part (2).
By Descent on Spaces, Lemma
\ref{spaces-descent-lemma-descending-property-finite-type}
we see that $X \to U$ is of finite type.
In the case of $\Spacesstack'_{fp, flat, proper}$
one additionally uses
Descent on Spaces, Lemma
\ref{spaces-descent-lemma-descending-property-finite-presentation},
\ref{spaces-descent-lemma-descending-property-flat}, and
\ref{spaces-descent-lemma-descending-property-proper}
in the last step.
\end{proof}

\end{document}
